\section{Erfüllung der Forschungsfragen}

Die erste Forschungsfrage nach einem flexiblen Rollen- und Berechtigungssystem kann nur teilweise als erfüllt bezeichnet werden. Die benutzerbezogene Zugriffskontrolle über Microsoft Entra ID wurde technisch solide umgesetzt, jedoch ist eine feingranulare rollenbasierte Differenzierung mit mehreren fachlichen Rollen, wie sie ursprünglich im Pflichtenheft geplant war, lediglich vorbereitet und noch nicht vollständig realisiert. Die theoretischen Grundlagen zu RBAC und ABAC hätten eine weitergehende Implementierung durchaus ermöglicht.

Die zweite Forschungsfrage nach einer skalierbaren und plattformunabhängigen Architektur ist als weitgehend beantwortet zu bewerten. Durch die serviceorientierte Aufteilung in vier klar getrennte Dienste, die Nutzung cloud-nativer Komponenten und die Docker-Containerisierung konnten die gestellten Anforderungen in der Praxis erfüllt werden. Die im theoretischen Kapitel erarbeiteten Skalierbarkeitsstrategien spiegeln sich direkt in der Wahl der Azure-Dienste wider.

Die dritte Forschungsfrage zur Erkennungsgenauigkeit von KI-Modellen bei der Dokumenttypenerkennung kann als beantwortet betrachtet werden. Die Entscheidung für GPT-4o wurde durch die theoretische Auseinandersetzung mit den Modellfamilien gut begründet, insbesondere hinsichtlich der Eignung für Zero-Shot-Szenarien und der einfachen Erweiterbarkeit auf neue Dokumenttypen. Die implementierte Nachvalidierung der KI-Ergebnisse sowie die Möglichkeit zur Benutzerkorrektur vor dem finalen Commit zeigen ein bewusstes Bewusstsein für die natürlichen Grenzen automatisierter Klassifikation und wurden gezielt als Absicherung in den Upload-Prozess integriert.

Die vierte Forschungsfrage nach der optimalen Such- und Filterarchitektur ist von allen Forschungsfragen am vollständigsten beantwortet. Die Implementierung folgt direkt dem theoretisch erarbeiteten dreischichtigen Aufbau aus Retrieval, Filterung und Facettenaggregation. Für den vorliegenden Anwendungsfall, bei dem Dokumente primär anhand strukturierter Metadaten wie Dokumenttyp, Jahr und Partner gesucht werden, stellt die lexikalische Volltextsuche über Azure AI Search eine vollständig ausreichende und praxisgerechte Lösung dar.